\scanpage{37}
% archive foliation: 36
\[
\unsure{\mathcal{S}|S}^\circ\longrightarrow\underline{\operatorname{Hom}}_{\unsure{k|\mathbf Z}\text{-}add}(A,\mathrm{Ab})
\]
Nous ferons cette étude dans le
cas où $\unsure{\mathcal{S}|S}=\widehat\Delta=\unsure{\mathrm{S}_s|\mathrm{Ss}}$, et \unsure{donc|nous}
\uline{\unsure{choisirons pour $C$ les suivants}}.

$1^\circ)$ $C$ est formée des groupes ab. $\frac12$ simpli-
ciaux \deleted{\unsure{libres de rangs finis}}\inserted{str. $\partial$-parfaits}, $A$ est la
sous-catégorie formée des
$\Phi(i,n)$ ($n\in\mathbf N^*$, \deleted{\unsure{j\geqslant0}} $i\geqslant1$) et $\Lambda^j\Psi_*$ ($j\geqslant0$).
\margin{$C$ abélienne}

On voit que \inserted{pour $X\in\operatorname{Ob}\widehat\Delta$ fixé} la connaissance des
$\operatorname{Hom}_{\unsure{\widehat\Delta|\Delta}}(X_*,M_*)$ comme foncteur
covariant en $M_*$ se réduit à
celle de $\operatorname{Hom}_{\unsure{\widehat\Delta|\Delta}}(X_*,\deleted{\unsure{L}}\varphi)$ pour
$\varphi$ dans $A$. On peut encore
réduire \deleted{\unsure{à}}\inserted{à prendre} $\varphi$ dans $A_0$ défini
ci-dessous, et même pour $\varphi$ prenant
les $\Phi^i$.

$2^\circ)$ $C_0$ \inserted{$\subset C$} est formée des groupes ab.
$\frac12$ simpliciaux \deleted{\unsure{str.}} spéciaux
(\deleted{donc}\inserted{i.e.} $\in C$ et les $\pi_i(L_*)$ sont
sans torsion), $A_0\subset C_0$ sous-cat.
de ceux qui sont de la forme
\[
\Phi^i=\Lambda^i\Phi_*\deleted{\unsure{\Psi}}\quad(i\geqslant1)
\quad\text{ou}\quad\Psi^j=\Lambda^j\Psi_*\quad(j\geqslant0).
\]
\margin{$C_0$ est additive mais non abélienne.}
